\documentclass[11pt]{article}
\usepackage[margin=1in]{geometry}
\usepackage{amsmath,amssymb}
\usepackage{booktabs}
\newcommand{\sbx}{S\nobreakdash-box}
\newcommand{\sbxes}{S\nobreakdash-boxes}

\title{Krakken-2048: Experimental Verification and Distinguisher Battery}
\author{}
\date{}

\begin{document}
\maketitle

\begin{abstract}
This note is a companion to the active-\sbx\ bound analysis of Krakken-2048. It
records two categories of result obtained by direct computation against the
reference implementation: (i) \emph{exact proofs} of the cryptographic
foundations the bound analysis relies on (the MDS branch number, the \sbx's\
differential and linear properties, and the bijectivity of the word-permutation
layers), each established by complete enumeration; and (ii) an \emph{empirical
distinguisher battery} (avalanche, rotational and RX-difference,
integral/saturation, cube / algebraic degree, one- and two-sided zero-sum, exact
algebraic-degree growth, differential clustering, impossible differentials,
boomerang / BCT, differential-linear, and invariant-subspace / \sbx-coset
structure) run against the compiled
permutation. The foundations are proven. The distinguisher battery finds no
distinguisher surviving the full eight rounds in any tested family. We are
explicit throughout that the battery is empirical and necessary-but-not-sufficient:
each test rules out a \emph{detectable} weakness of its kind at the sampled
scale, not the existence of a more subtle attack, and several have stronger
(MILP-based) forms that remain future work.
\end{abstract}

\section{Exact proofs of the foundations}

The following properties underlie every active-\sbx\ bound in the companion
analysis. They had previously been inherited by analogy from AES/Whirlpool; here
each is verified exhaustively for the constants as implemented.

\begin{table}[h]
\centering
\begin{tabular}{lll}
  \toprule
  \textbf{Property} & \textbf{Method} & \textbf{Result} \\
  \midrule
  MDS branch number & all $12{,}869$ square submatrices over GF($2^8$) & $9$ (maximal) \\
  \sbx\ bijection & full table check & permutation of $0\!-\!255$ \\
  \sbx\ diff.\ uniformity & full $256\times256$ DDT & $4 \Rightarrow \mathrm{DP}=2^{-6}$ \\
  \sbx\ linearity & full $256\times256$ LAT & bias $16 \Rightarrow \mathrm{corr}^2=2^{-6}$ \\
  \sbx\ boomerang uniformity & full $256\times256$ BCT & $6 \Rightarrow$ switch $2^{-5.42}$ \\
  $\pi$ word permutation & index-map bijection check & bijection \\
  InkCloud word permutation & index-map bijection check & bijection \\
  InkCloud multiplier & $\gcd(7,32)=1$ & invertible, $7^{-1}=23 \bmod 32$ \\
  \bottomrule
\end{tabular}
\caption{Exhaustively-proven foundational properties (reduction polynomial
$\mathtt{0x11D}$; MDS row $(1,1,4,1,8,5,2,9)$; ABYSSAL \sbx).}
\end{table}

These establish that the per-\sbx\ weight is exactly $2^{-6}$ for both
differential and linear cryptanalysis, that the diffusion matrix achieves the
maximal branch number $9$, and that the permutation layers preserve the full
state (are bijective). The inverse permutation was additionally implemented and
verified to round-trip bit-exactly with the forward permutation in both
directions at all eight round counts.

\section{Empirical distinguisher battery}

Each test below was run against the compiled reference permutation. For each we
state what a positive (weak) result would look like and what was observed.

\subsection{Avalanche (strict diffusion)}
Flipping a single input bit should change about half of the $2048$ output bits.
Measured over $200$ trials per round:

\begin{table}[h]
\centering
\begin{tabular}{lcccc}
  \toprule
  \textbf{Round} & \textbf{mean} & \textbf{min} & \textbf{max} & \textbf{std} \\
  \midrule
  1 & 0.4996 & 0.4673 & 0.5332 & 0.0114 \\
  2 & 0.4998 & 0.4707 & 0.5454 & 0.0119 \\
  3 & 0.4990 & 0.4639 & 0.5283 & 0.0107 \\
  4 & 0.5004 & 0.4634 & 0.5322 & 0.0115 \\
  5 & 0.5009 & 0.4604 & 0.5283 & 0.0112 \\
  6 & 0.5004 & 0.4722 & 0.5386 & 0.0108 \\
  7 & 0.4994 & 0.4712 & 0.5312 & 0.0113 \\
  8 & 0.4997 & 0.4717 & 0.5337 & 0.0119 \\
  \bottomrule
\end{tabular}
\caption{Fraction of output bits changed by a one-bit input flip. Full avalanche
(mean within $1\%$ of $0.5$) is reached at round $1$ and held through round $8$.}
\end{table}

\subsection{Strict avalanche criterion, completeness, and monobit balance}
The avalanche test above reports only the \emph{mean} fraction of output bits
flipped. A mean of exactly $0.5$ is, however, compatible with a badly structured
per-bit dependency: an input bit that always flips a given output bit, and never
flips another, averages to $0.5$ while violating diffusion. To see this we
measure the full strict-avalanche-criterion (SAC) matrix (Webster and Tavares, 1986): for
every pair (input bit $i$, output bit $j$),
\[
  \mathrm{SAC}[i][j] \;=\; \Pr\!\big[\text{output bit } j \text{ flips}
  \;\big|\; \text{input bit } i \text{ is flipped}\big]
  \qquad(\text{ideal } \tfrac12),
\]
estimated over $N=4096$ random base states per input bit ($2048\times2048=
4.19\times10^{6}$ cells, $2\times2048\times N$ permutation calls per round). From
the same data we extract two derived properties: \emph{completeness} (every cell
strictly between $0$ and $N$, i.e.\ every input bit influences every output bit
but none deterministically) and a per-output-bit \emph{monobit} frequency over
random inputs (ideal $\tfrac12$). At $N=4096$ the two-sided $5\sigma$ threshold
on a single proportion is $|p-\tfrac12|>0.0391$.

\begin{table}[h]
\centering
\begin{tabular}{lccccc}
  \toprule
  \textbf{Round} & \textbf{SAC max} & \textbf{SAC mean} &
  \textbf{cells $>5\sigma$} & \textbf{dead / stuck} & \textbf{monobit max} \\
   & $|p{-}\tfrac12|$ & $|p{-}\tfrac12|$ & (exp.\ $\approx2.4$) & cells & $|f{-}\tfrac12|$ \\
  \midrule
  1 & 0.0415 & 0.0062 & 4 & 0 / 0 & 0.0276 \\
  2 & 0.0405 & 0.0062 & 2 & 0 / 0 & 0.0295 \\
  3 & 0.0422 & 0.0062 & 5 & 0 / 0 & 0.0305 \\
  4 & 0.0422 & 0.0062 & 2 & 0 / 0 & 0.0249 \\
  5 & 0.0413 & 0.0062 & 3 & 0 / 0 & 0.0264 \\
  6 & 0.0427 & 0.0062 & 1 & 0 / 0 & 0.0293 \\
  7 & 0.0435 & 0.0062 & 4 & 0 / 0 & 0.0271 \\
  8 & 0.0413 & 0.0062 & 4 & 0 / 0 & 0.0303 \\
  \bottomrule
\end{tabular}
\caption{SAC matrix, completeness, and monobit balance ($N=4096$ per input bit,
$5\sigma=0.0391$). ``SAC mean $|p{-}\tfrac12|$'' is the mean absolute deviation
across all $4.19\times10^{6}$ cells; the value $0.0062$ is exactly that expected
of an ideal $\mathrm{Binomial}(N,\tfrac12)$ proportion
($\approx0.399\,\sigma$). ``cells $>5\sigma$'' tracks the random-expected $2.4$
false positives. ``dead'' cells never flip (incomplete diffusion), ``stuck''
cells always flip (deterministic coupling); both are zero at every round.}
\end{table}

The SAC matrix is statistically indistinguishable from ideal at \emph{every}
round including round~$1$: the mean cell deviation is exactly the random-sampling
value, and the one-to-five cells exceeding $5\sigma$ are consistent with the
$\approx2.4$ false positives expected over $4.19\times10^{6}$ cells. The
completeness count is decisive: with zero dead and zero stuck cells at all
rounds, every one of the $2048$ input bits provably influences every one of the
$2048$ output bits, and none does so deterministically---the per-bit form of the
full-state first-round diffusion claimed by the mean-avalanche test. No
output-bit monobit bias was detected at any round. This is a stronger negative
result than the mean avalanche: a permutation with a correct mean but a
structured dependency matrix would fail here, and none of the $2048\times2048$
couplings deviates beyond sampling noise.

\subsection{NIST SP 800-22 statistical test suite on the sponge keystream}
The probes above target the bare permutation. We additionally subjected the full
sponge \emph{keystream}---the permutation as used in the stream-cipher mode of
the Krakken-Disk volume encryptor---to the NIST Statistical Test Suite
(SP~800-22). Three independent $5500$\,MB encrypted volume containers were
produced. Each container is initialised \emph{empty}, so its plaintext is
essentially the all-zero sector image; the encryptor forms
$\mathit{ciphertext}_i = \mathit{plaintext}_i \oplus
\mathrm{Sponge}(\mathit{key}\,\|\,\mathit{nonce}\,\|\,\mathrm{LE64}(i))$,
so up to the per-sector random nonces ($24$ bytes per $4096$-byte sector) and a
small fixed header the ciphertext \emph{is} the Krakken sponge keystream. Testing
the ciphertext of an empty container is therefore, in effect, testing the raw
keystream, with no plaintext entropy masking the permutation.

Each file was assessed over $10$ sequences of $10^{6}$ bits (the standard STS
configuration). All fifteen test families---Frequency, Block Frequency,
Cumulative Sums, Runs, Longest Run, Binary Matrix Rank, Spectral (FFT),
Non-Overlapping Template ($148$ templates), Overlapping Template, Maurer's
Universal, Approximate Entropy, Random Excursions and its Variant, Serial, and
Linear Complexity---passed for all three files: every test's proportion of
passing sequences met or exceeded the NIST minimum ($8/10$, or $5/6$ for the
random-excursion tests), and every uniformity $P$-value was non-degenerate.%
\footnote{The suite overwrites its summary report on each run; the complete
fifteen-family report is retained for \texttt{random3.bin} (no failures) and the
other two files were logged at the same all-pass level.}

As an independent high-throughput cross-check, a custom monobit and runs test was
run over a larger $10^{8}$-bit window of each file (Table~\ref{tab:nist_ht}); all
$P$-values are unremarkable and every byte-level Shannon entropy exceeded
$7.999$ bits/byte.

\begin{table}[h]
\centering
\begin{tabular}{lccc}
  \toprule
  \textbf{File} & \textbf{Monobit $P$} & \textbf{Runs $P$} & \textbf{Entropy (bits/byte)} \\
  \midrule
  \texttt{random1.bin} & 0.7165 & 0.5258 & $>7.999$ \\
  \texttt{random2.bin} & 0.2247 & 0.3089 & $>7.999$ \\
  \texttt{random3.bin} & 0.9925 & 0.5349 & $>7.999$ \\
  \bottomrule
\end{tabular}
\caption{High-throughput monobit and runs cross-check over $10^{8}$ bits per
file. All $P$-values exceed the $0.01$ significance threshold.}
\label{tab:nist_ht}
\end{table}

This is a mode-level empirical randomness check on the sponge keystream: like the
rest of the battery it is necessary but not sufficient, ruling out detectable
first- and second-order statistical bias at the sampled scale
($\approx10^{7}$--$10^{8}$ bits) rather than establishing a security property of
any construction built on the permutation.

\subsection{Sponge-hash collisions and output distribution}
The NIST suite above tests the keystream; here we test the sponge \emph{hash}
($256$-bit digest, rate $160$ bytes) for collision and distribution behaviour, a
family the rest of this note does not cover. Three quick probes were run against
the compiled hash.

\textbf{Birthday collisions.} We hashed $N=10^{6}$ random messages, truncated
each digest to $t$ bits, and counted colliding pairs. For an ideal random
function the expected number of colliding pairs is
$\binom{N}{2}/2^{t}\approx N^{2}/2^{t+1}$; a structural collision weakness would
appear as a large excess at some truncation width.

\begin{table}[h]
\centering
\begin{tabular}{rrrr}
  \toprule
  \textbf{$t$ (bits)} & \textbf{expected} & \textbf{observed} & \textbf{obs/exp} \\
  \midrule
  28 & 1862.6 & 1835 & 0.985 \\
  30 &  465.7 &  495 & 1.063 \\
  32 &  116.4 &  122 & 1.048 \\
  34 &   29.1 &   27 & 0.928 \\
  36 &    7.28 &    5 & 0.687 \\
  38 &    1.82 &    2 & 1.100 \\
  40 &    0.46 &    1 & 2.199 \\
  \bottomrule
\end{tabular}
\caption{Truncated-digest birthday collisions, $N=10^{6}$ messages. Observed
counts track the ideal $N^{2}/2^{t+1}$ closely wherever the expectation is
statistically meaningful ($t\le34$); the apparent ratios at $t\ge36$ are
small-count Poisson fluctuations (expectation below $7$). No truncation width
shows a collision excess.}
\end{table}

\textbf{Output diffusion / near-collisions.} For $2\times10^{5}$ random message
pairs differing in a single input bit, the full $256$-bit digest difference had
mean Hamming weight $127.98$ (ideal $128$) and standard deviation $8.00$ (ideal
$\sqrt{256\cdot\tfrac14}=8$); the minimum over all pairs was $93$ and the maximum
$165$. The minimum staying far from $0$ rules out cheap one-bit near-collisions.

\textbf{Digest byte uniformity.} Over $2\times10^{5}$ digests
($6.4\times10^{6}$ output bytes), the byte-value frequencies gave
$\chi^{2}=273.8$ on $255$ degrees of freedom (ideal $255\pm22.6$; $z=+0.83$),
consistent with a uniform output distribution.

All three probes are consistent with an ideal random function; as elsewhere they
are empirical and sampling-bounded, excluding gross collision or distribution
defects at the $\sim\!10^{6}$-message scale rather than establishing collision
resistance analytically.

\subsection{Rotational}
Rotational cryptanalysis exploits that XOR, AND, and bit-rotation commute with
rotating the whole state. We compare $\mathrm{rotate}(P(X))$ against
$P(\mathrm{rotate}(X))$: a rotation-invariant design makes these equal, a sound
design makes them differ in $\sim\!1024$ of $2048$ bits (random). Measured mean
Hamming distance was $1022.9$–$1025.2$ across rounds $1$–$8$ for rotation
amount $1$, and $1022.9$–$1024.9$ at eight rounds across rotation amounts
$\{1,7,13,17,32,63\}$. No rotational relation was detected at any round; the
relation is destroyed from round $1$, consistent with the non-symmetric round
constants and the carry behaviour of the PRESSURE additions breaking rotational
invariance.

\subsection{Integral / saturation}
Saturating one input byte through all $256$ values and XOR-summing the outputs
should yield a balanced (zero-sum) signature in many output bytes if the design
has a surviving integral property. Random expectation is $\sim\!1$ zero byte of
$256$. Measured averages were $0.25$–$1.25$ zero bytes across rounds $1$–$8$
(samples at or below the random baseline at every round). No integral
distinguisher was found at any round---including round $1$, a stronger outcome
than AES, whose classical integral survives three rounds; Krakken's full-state
first-round diffusion removes the balanced property immediately.

\subsection{Cube / algebraic degree}
XOR-summing the output over all assignments of a $k$-bit input cube reveals the
top-degree ANF coefficient; a structural zero across many constant settings is a
cube distinguisher. Using $k=8$ and requiring the cube-sum bit to vanish across
$20$ random constants (chance baseline $2048\cdot2^{-20}\approx0.002$, so any
survivor is real signal), the count of consistent-zero output bits was:

\begin{table}[h]
\centering
\begin{tabular}{lc}
  \toprule
  \textbf{Round} & \textbf{consistent-zero bits (of 2048)} \\
  \midrule
  1 & 32 \\
  2 & 0 \\
  3 & 0 \\
  4 & 0 \\
  8 & 0 \\
  \bottomrule
\end{tabular}
\caption{Cube / algebraic-degree distinguisher. A genuine distinguisher shows a
count that decays to zero as algebraic degree builds; here it does so by round
$2$. The round-$1$ value ($32$) is a real but expected low-degree artifact of a
single round, far below the eight-round permutation.}
\end{table}

\subsection{Zero-sum (forward and two-sided)}
A zero-sum set has both input-XOR and output-XOR equal to zero. The
\emph{forward} test built affine input subspaces (input-XOR $=0$ by
construction, dimension $8$, $24$ subspaces) and found no round at which the
output-XOR was also zero ($0$ full zero-sums, $0$ consistent-zero bits at every
round). The \emph{two-sided} test (using the verified inverse permutation,
dimension $9$, $12$ subspaces) placed a subspace at every split
$r_{\text{back}}+r_{\text{fwd}}=8$ and required input-XOR and output-XOR to
vanish simultaneously. Only the trivial edge splits ($0\!+\!8$ and $8\!+\!0$, in
which one side is the raw subspace) showed any zero-sum; the degree saturates
within two rounds from each side, and no split with $r_{\text{back}}\ge1$ and
$r_{\text{fwd}}\ge1$ produced a two-sided zero-sum. No zero-sum distinguisher was
found over the full eight rounds.

\subsection{Exact algebraic degree}
This is the test most capable of revealing a higher-order differential or cube
weakness: rather than checking a fixed cube, it computes the \emph{exact}
algebraic-normal-form degree of each output bit as a Boolean function of a chosen
set of $n=10$ input variables (via the M\"obius transform over all $2^{10}$
assignments, other inputs fixed). A degree that plateaus below the maximum $n$ as
rounds increase would be a distinguisher; a healthy permutation saturates to
near-maximal degree within a few rounds. We tested six independent variable
placements and report the weakest (minimum) output-bit degree, since an attacker
chooses the variables that minimise it.

\begin{table}[h]
\centering
\begin{tabular}{lccl}
  \toprule
  \textbf{Round} & \textbf{max degree} & \textbf{min degree (weakest bit)} & \textbf{} \\
  \midrule
  1 & 10 & 5 & degree still building \\
  2 & 10 & 8 & saturated \\
  3 & 10 & 8 & saturated \\
  4 & 10 & 8 & saturated \\
  8 & 10 & 8 & saturated \\
  \bottomrule
\end{tabular}
\caption{Exact algebraic degree over $n=10$ input variables (maximum possible
degree $10$). Degree reaches near-maximal by round $2$ and holds.}
\end{table}

The minimum degree of $8$ from round $2$ onward warranted scrutiny, since a
plateau below maximal is the signature of a distinguisher. Direct inspection of
the full degree distribution at eight rounds resolves it as benign: of the $2048$
output bits, $1027$ have degree $10$, $1020$ have degree $9$, and exactly
\emph{one} has degree $8$ (and that bit was verified to depend on all $10$
variables, so its degree is not capped by missing dependence). This is precisely
the distribution expected of $2048$ near-random degree-$\le 10$ Boolean functions:
the top-degree ANF coefficient is present or absent with probability $\approx
1/2$, so a small number of bits randomly fall to degree $9$ and, rarely, $8$. A
genuine weakness would instead show many output bits stuck at low degree with the
count persisting or growing across rounds; the observed behaviour is the opposite
(full saturation by round $2$ with a single statistical straggler). No
low-degree distinguisher is present over the tested variable sets.

\subsection{Division-property integral (bit-based, MILP)}
\label{sec:divprop}
The integral test above (byte saturation) is a sampling test; its analytical
counterpart is the \emph{bit-based division property} (two-subset / conventional
CBDP, Todo--Xiang et al.), which propagates a division trail through the round
function as a mixed-integer program and proves an output bit \emph{balanced}
(zero-sum over the chosen input cube) exactly when its division trail is
infeasible. Unlike the empirical battery, a bit certified balanced this way is a
\emph{proof}, not a statistical observation. This is the analysis previously
listed as future work; we report it here as a working, sound harness together
with its current results.

\paragraph{Model.} One Krakken round is modelled bit-for-bit: the three
GF(2)-linear layers (the \textsc{MDS}-based diffusion, the butterfly diffusion,
and the \textsc{InkCloud} word-shuffle) by their exact matrices extracted from
the reference implementation, the ARX \textsc{Pressure} layer by a sound
full-adder carry-chain over-approximation, and each ABYSSAL \sbx\ by its
division-property transitions. The circuits are validated bit-for-bit against the
compiled permutation and the ARX division rule against exhaustive enumeration of a
$4$-bit adder, so the model is \emph{sound}: every bit it certifies balanced is a
genuine zero-sum bit, and the conservative directions (the ARX over-approximation
and the inequality \sbx\ model below) can only cause it to \emph{miss} balanced
bits, never to invent one.

\paragraph{Making it tractable.} The textbook ``selector'' \sbx\ model encodes
the $\approx\!2015$ minimal division trails of the $8$-bit \sbx\ as that many
binary selector variables \emph{per \sbx\ instance}---about $5\times10^{5}$
binaries per round across the $256$ \sbxes\ of the nonlinear layer. At that size
the solver (SCIP) could not resolve a \emph{single} output bit, even without a
time limit (the one-round model has $937{,}952$ variables; every output bit hit
the per-bit time cap). We replace it with a compact \emph{inequality} model: a set
of linear inequalities over the $16$ input/output bits with \emph{no} auxiliary
variables, generated as one ``no-good'' exclusion cut per impossible division
transition and then greedily merged. The result is $661$ inequalities and is
\emph{exact}---verified over all $2^{16}$ points of the \sbx\ division polytope to
retain every one of the $62{,}741$ valid transitions and exclude all $2{,}795$
impossible ones, reproducing the trail table and the algebraic degree $7$. This
cuts the one-round model from $937{,}952$ to $422{,}112$ variables and is what
makes the program solvable at all.

\paragraph{Current result.} With the compact model the round-$1$ program solves
to optimality---for a saturated input word as the cube, output bit~$0$ is proved
\emph{reachable} (i.e.\ not balanced) in $\approx\!160$\,s, whereas the selector
model never terminated on any bit. This is consistent with the empirical integral
finding above: full-state first-round diffusion already removes byte balance, so
no surviving integral is expected even at one round. The exhaustive
determination---sweeping all $2048$ output bits and increasing the round count
until no balanced bit remains, which would fix the integral-distinguisher
length---is a long-running computation; the operational harness and its round-$1$
confirmation are reported here, with the multi-round sweep to follow. The harness
is released alongside this note so the result can be reproduced and extended.

\subsection{Differential clustering (empirical differential probability)}
The active-\sbx\ bound in the companion analysis bounds the best single
\emph{characteristic}; a differential attack instead exploits a \emph{differential}
$(\delta_{\text{in}}\!\to\!\delta_{\text{out}})$, summing the probabilities of all
characteristics sharing that input/output pair (the cluster). A differential can
therefore have higher probability than any single characteristic. We probe this
empirically: fix the minimal-trail input difference identified in the companion
analysis ($4$ active bytes, all in lane~$7$, in words $0,12,19,28$), push
$200{,}000$ random message pairs $(x, x\oplus\delta_{\text{in}})$ through $r$
rounds of the compiled permutation, and tally the resulting output differences on
a $16$-byte ($128$-bit) projection. A surviving cluster would appear as some
projected output difference recurring far more often than chance; at this scale
($2\times10^5$ trials over $2^{128}$ cells) the chance baseline is essentially
zero, so any projected difference recurring even twice is candidate signal
(detection floor $\approx 2^{-17}$).

\begin{table}[h]
\centering
\begin{tabular}{lccl}
  \toprule
  \textbf{Round} & \textbf{distinct diffs} & \textbf{max recurrence} & \textbf{} \\
  \midrule
  1 & 200{,}000 & 1 & flat \\
  2 & 200{,}000 & 1 & flat \\
  3 & 200{,}000 & 1 & flat \\
  4 & 200{,}000 & 1 & flat \\
  5 & 200{,}000 & 1 & flat \\
  6 & 200{,}000 & 1 & flat \\
  7 & 200{,}000 & 1 & flat \\
  8 & 200{,}000 & 1 & flat \\
  \bottomrule
\end{tabular}
\caption{Differential clustering on the minimal-trail input difference,
$200{,}000$ message pairs per round, $128$-bit output-difference projection. Every
projected output difference was unique at every round (maximum recurrence $1$
against a chance threshold of $2$); no differential cluster was detected. Detection
floor $\approx 2^{-17}$.}
\end{table}

No differential cluster with probability above the $\approx 2^{-17}$ detection
floor was found for the minimal-trail input difference at any round from one to
eight. This is direct empirical evidence against a fat differential cluster on the
trail the MILP bound rides; it is sampling-bounded and does not exclude a fainter
cluster (the proven single-characteristic bound is far below any reachable
sampling sensitivity), so it complements rather than replaces the analytical
single-trail bound.

\subsection{RX-difference (rotational-XOR) bias}
The plain rotational test above rules out exact rotational symmetry; a finer
probe is whether a \emph{rotational-XOR} (RX) difference --- a fixed difference
imposed between a state and a rotated copy --- propagates with detectable bias.
We measured, over $500{,}000$ trials per configuration at one round, the maximum
per-output-bit bias across all $2048$ output bits (the $5\sigma$ detection
threshold at this sample size is $0.003536$). A single-word sweep over rotation
constants $\{7,13,17,19,31,37\}$ and injection words $\{0,8,16\}$ left $17$ of
$18$ configurations clearly below threshold (range $0.002228$--$0.002796$); the
lone marginal exceedance (\texttt{rx\_rot=17}, \texttt{inj\_word=8}: bias
$0.003572$, one bit of $2048$, $\approx\!1.01\times$ threshold) is consistent
with the expected false-positive rate. A heavy-mode sweep over all $63$
non-trivial rotation constants for four multi-word injection blocks ($252$
configurations) failed to reproduce it: every configuration stayed below
threshold, the highest bias observed being $0.003472$ ($1.8\%$ below the
$5\sigma$ threshold). No RX-difference distinguisher was detected.

\subsection{Impossible differentials}
For each of four input-bit positions ($0,64,512,1984$) and four output-byte
positions ($0,31,127,255$), $200{,}000$ random pairs $(x,x\oplus\delta)$ were
pushed through the reduced-round permutation and the output-difference histogram
at the chosen byte was tallied; any non-trivial output difference with count
zero is an impossible-differential candidate. At both one and two rounds every
completed configuration was clean ($0/255$ non-trivial differences missing). No
impossible-differential candidate was found.

\subsection{Boomerang / BCT}
The full $256\times256$ Boomerang Connectivity Table of the \sbx\ was computed by
enumeration. Its maximum entry (excluding the trivial $(a,0)$ and $(0,b)$ rows)
is $6$, at $(a,b)=(\mathtt{0x01},\mathtt{0x05})$, with non-zero entries
distributed as $6$ at $510$ pairs, $4$ at $255$ pairs, $2$ at $31{,}620$ pairs.
A BCT maximum of $6$ exceeds the DDT-flat ideal of $2$, giving a per-\sbx\
boomerang switch probability of $2^{-5.42}$ --- a property of the \sbx\ alone and
round-independent. An empirical $50{,}000$-trial boomerang quartet experiment
with input difference $\alpha=\mathtt{0x01}$ at byte~$0$ gave a maximum
output-difference count of $232$ at one round (uniform expectation $195.3$, ratio
$1.19$) and $229$ at two rounds (ratio $1.17$), both well below the $2.0$ anomaly
threshold. The elevated \sbx\ BCT entry does not propagate to a full-permutation
boomerang distinguisher at any tested round count.

\subsection{Differential-linear}
This test splits the permutation into a top and bottom half and so requires at
least two rounds. At two rounds, five $(r_0,\delta,\gamma)$ configurations were
probed over $200{,}000$ trials each (the $5\sigma$ bias threshold at this sample
size is $0.011180$). The highest observed correlation was $|\mathrm{corr}|=
0.003380$ --- less than a third of the threshold --- and all five configurations
were below it. No differential-linear bias was detected.

\subsection{Invariant subspace / \sbx-coset structure}
An exhaustive enumeration of all $(a,g)$ pairs satisfying $S(a)\oplus S(a\oplus
g)=g$ found $282$ dimension-one coset hits across $141$ distinct generators (of
$255$) --- the candidate seeds of an invariant-subspace trail. None of the $282$
survived a full round in either the constant-base probe or a random-base
empirical probe ($2{,}000$ trials per sampled pair), at one or two rounds. An
exhaustive check of all $32{,}385$ dimension-two subspaces found no coset hit at
all. The dimension-one hit count is unremarkable for an $8$-bit \sbx; the
decisive observation is that the structure does not survive the round function,
so it cannot seed an invariant-subspace trail.

\subsection{Structural sanity checks}
Independently of the above, the compiled permutation passed a battery of gross
structural checks (sampled, not proofs): bijectivity on a sample
($2000/2000$ distinct outputs), absence of fixed points ($0$ among $503$ states,
including the all-zero and all-one states), and a sampled subspace-trail collapse
test (random low-dimensional affine subspaces show no rank collapse under the
permutation). All passed. These rule out several gross structural weaknesses but
do not substitute for an analytical invariant-subspace / subspace-trail study.

\section{Scope and honest limitations}

The Section~2 foundations are exact proofs. The Section~3 battery is
\emph{empirical}: each test rules out a detectable distinguisher of its kind at
the sampled scale, and passing is necessary but not sufficient for security.
Specifically:
\begin{itemize}\itemsep2pt
  \item The rotational test is an empirical distinguisher, not a proof of
        rotational-XOR security; an RX-trail search with explicit carry
        probabilities (MILP) is the stronger form and remains future work.
  \item The integral test in this section is the classical (byte-saturation)
        form. Its analytical counterpart, the bit-based division-property integral
        (two-subset CBDP, MILP), is now implemented (Section~\ref{sec:divprop}): the harness is
        operational and sound, the round-1 program solves to optimality, and the
        exhaustive multi-round sweep that would fix the integral-distinguisher
        length is an ongoing computation rather than future work.
  \item The cube and zero-sum tests use fixed cube/subspace dimensions; they are
        not exhaustive ANF-degree bounds.
  \item The algebraic-degree test samples six variable placements; an attacker
        chooses the placement freely, so a carefully-positioned higher-order
        distinguisher is not excluded. The exhaustive form is a division-property
        degree bound (MILP) over all variable choices, which remains future work
        and is the highest-value remaining self-conducted analysis.
  \item The impossible-differential, boomerang, differential-linear, and
        invariant-subspace probes were run only at one and two rounds (the
        regime where a structural distinguisher, if present, is cheapest to
        detect); they bound reduced-round behaviour, not the full eight rounds,
        and the impossible-differential runs were additionally truncated by
        their time limit before every input-bit position completed.
  \item Algebraic (Gr\"obner-basis) attacks, the interaction of the permutation
        with the sponge construction, and---most importantly---adversarial
        review by independent cryptanalysts are not addressed here.
\end{itemize}
Taken together with the proven active-\sbx\ bounds in the companion analysis,
these results constitute a thorough self-conducted first-pass evaluation: the
foundations are proven, and no distinguisher of the tested families survives the
full permutation. They are offered as evidence supporting the design, not as a
substitute for external cryptanalysis.

\end{document}
